%=============================================================
% Silabus Mata Kuliah: Pemrograman Game (Berbasis OBE)
%=============================================================
\documentclass[12pt,a4paper]{article}
\usepackage[utf8]{inputenc}
\usepackage[T1]{fontenc}
\usepackage[indonesian]{babel}
\usepackage{geometry}
\usepackage{longtable}
\usepackage{array}
\usepackage{booktabs}
\usepackage{enumitem}
\usepackage[hidelinks,breaklinks]{hyperref}
\geometry{margin=2.5cm}

\begin{document}

\begin{center}
\textbf{\Large RENCANA PEMBELAJARAN SEMESTER (RPS)}\\
\textbf{Berbasis Outcome-Based Education (OBE)}\\[0.5cm]
\textbf{PROGRAM STUDI TEKNIK INFORMATIKA}\\
\textbf{FAKULTAS TEKNIK}\\
\textbf{UNIVERSITAS LOREM IPSUM}\\
\textbf{Mata Kuliah: Pemrograman Game}\\
\textbf{Tahun Ajaran: 2024/2025}
\end{center}

\vspace{1cm}

% ============================================================
% 1. IDENTITAS MATA KULIAH
% ============================================================
\section*{1. Identitas Mata Kuliah}
\begin{tabular}{ll}
Nama Program Studi   & : Teknik Informatika \\
Nama Mata Kuliah     & : Pemrograman Game \\
Kode Mata Kuliah     & : TIF-4xx \\
Semester            & : 5 (Lima) \\
SKS / Bobot Kredit   & : 3 SKS (2 Teori, 1 Praktikum) \\
Dosen Pengampu       & : [Nama Dosen], M.Kom. \\
Konsultasi          & : [Jadwal Konsultasi] \\
Email               & : [email@universitas.ac.id] \\
Platform Pembelajaran & : Unity Learn, GitHub, Discord \\
Tanggal Penyusunan   & : \today \\
Versi               & : 2.0 (Update 2024)
\end{tabular}

\vspace{0.5cm}

% ============================================================
% 2. CAPAIAN PEMBELAJARAN LULUSAN (CPL)
% ============================================================
\section*{2. Capaian Pembelajaran Lulusan (CPL)}
CPL yang dibebankan pada mata kuliah ini mencakup kompetensi lulusan dalam empat dimensi, mengacu pada standar industri game development 2024:
\begin{itemize}[leftmargin=*]
  \item \textbf{CPL‑1 (Pengetahuan)} : Menguasai konsep dasar pengembangan game, alur kerja Unity Engine, prinsip‑prinsip pemrograman C\#, serta game design patterns yang relevan.
  \item \textbf{CPL‑2 (Keterampilan Umum)} : Mampu menerapkan pemikiran logis, kritis, dan sistematis dalam merancang serta mengembangkan solusi game dengan pendekatan iterative development.
  \item \textbf{CPL‑3 (Keterampilan Khusus)} : Mampu mengoperasikan Unity Editor (versi 2024+), menulis skrip C\# yang efisien, menerapkan teknik AI, fisika, animasi, UI/UX, audio, dan optimization dalam konteks game development.
  \item \textbf{CPL‑4 (Sikap)} : Menunjukkan sikap profesional, kolaboratif version control (Git), dan etis dalam tim pengembangan game sesuai industry best practices.
\end{itemize}

\vspace{0.5cm}

% ============================================================
% 3. CAPAIAN PEMBELAJARAN MATA KULIAH (CPMK)
% ============================================================
\section*{3. Capaian Pembelajaran Mata Kuliah (CPMK)}
Setelah menyelesaikan mata kuliah ini, mahasiswa diharapkan mampu:
\begin{enumerate}[label=CPMK\arabic*:, leftmargin=*]
  \item Menganalisis konsep dasar game development, game design, genre game, dan menjelaskan workflow pengembangan game menggunakan Unity Engine 2024+.
  \item Menerapkan pemrograman C\# untuk Unity dengan menguasai struktur skrip, fungsi dasar (\texttt{Start()}, \texttt{Update()}, \texttt{Awake()}, \texttt{FixedUpdate()}), variabel, kontrol alur, konsep OOP, dan modern C\# features.
  \item Mengoperasikan Unity Editor secara profesional dalam mengelola GameObject, komponen (Transform, Rigidbody, Collider), prefab, scene, asset management, dan Package Manager.
  \item Mengimplementasikan sistem kontrol pemain menggunakan Unity Input System (versi baru), menggerakkan karakter, dan menangani interaksi pemain dengan game world.
  \item Mengaplikasikan sistem fisika Unity (Rigidbody, Collider, Physics Material, Force, Torque) untuk simulasi gerak, gravitasi, dan tumbukan yang realistis.
  \item Membuat animasi karakter dan objek menggunakan Animator Controller, Animation Layers, State Machines, Blend Trees, dan mengintegrasikan animasi dengan gameplay.
  \item Merancang dan mengimplementasikan UI/UX game yang efektif meliputi HUD, menu navigasi, sistem scoring, feedback visual, dan Unity UI Toolkit.
  \item Mengintegrasikan audio (sound effect dan background music) menggunakan AudioSource, AudioListener, Audio Mixer, dan spatial audio.
  \item Merancang level design dan game mechanics yang menarik serta menyeimbangkan progresi, tantangan, reward, dan player engagement.
  \item Mengimplementasikan AI sederhana untuk NPC menggunakan NavMesh, State Machines, Behavior Trees, dan basic pathfinding algorithms.
  \item Menganalisis serta mengoptimalkan performa game melalui Unity Profiler, debugging, memory management, dan teknik optimasi aset serta kode.
  \item Membangun dan mendeploy game ke berbagai platform (PC, Mobile, Web, Console) dengan memahami proses build, settings, dan distribusi.
  \item Mengembangkan proyek game lengkap yang mencakup seluruh aspek pembelajaran dari konsep hingga publikasi dengan version control dan documentation.
\end{enumerate}

\vspace{0.5cm}

% ============================================================
% 4. Sub‑CPMK / Indikator Pencapaian
% ============================================================
\section*{4. Sub‑CPMK / Indikator Pencapaian}
Berikut adalah indikator terukur untuk masing‑masing CPMK:
\begin{enumerate}[label=Sub‑CPMK\arabic*:, leftmargin=*]
  \item \textbf{Analisis Konsep Game Development}: Mampu menjelaskan game loop, genre classification, dan Unity development workflow dengan diagram alir (CPMK1).
  \item \textbf{C\# Scripting Fundamentals}: Mampu menulis script C\# dengan proper structure, variable declarations, dan basic functions (CPMK2).
  \item \textbf{Unity Editor Operations}: Mampu mengelola GameObject hierarchy, component configuration, dan scene organization (CPMK3).
  \item \textbf{Player Control Implementation}: Mampu mengimplementasikan character movement, input handling, dan state management (CPMK4).
  \item \textbf{Physics Simulation}: Mampu menerapkan Rigidbody dynamics, collision detection, dan physics materials (CPMK5).
  \item \textbf{Animation System}: Mampu membuat Animator Controller, state transitions, dan blend trees (CPMK6).
  \item \textbf{UI/UX Development}: Mampu merancang HUD, menu systems, dan interactive UI elements (CPMK7).
  \item \textbf{Audio Integration}: Mampu mengimplementasikan sound effects, background music, dan spatial audio (CPMK8).
  \item \textbf{Level Design Principles}: Mampu merancang level progression, difficulty curves, dan player engagement mechanics (CPMK9).
  \item \textbf{AI Implementation}: Mampu membuat NPC behavior dengan NavMesh navigation dan basic state machines (CPMK10).
  \item \textbf{Performance Optimization}: Mampu menggunakan Unity Profiler, identify bottlenecks, dan implement optimizations (CPMK11).
  \item \textbf{Build and Deployment}: Mampu configure build settings, optimize untuk target platforms, dan create deployment packages (CPMK12).
  \item \textbf{Project Integration}: Mampu mengembangkan complete game project dengan proper documentation dan portfolio presentation (CPMK13).
\end{enumerate}

\vspace{0.5cm}

% ============================================================
% 5. Materi Pembelajaran (Bahan Kajian)
% ============================================================
\section*{5. Materi Pembelajaran (Bahan Kajian)}
\begin{enumerate}[leftmargin=*]
  \item Konsep Dasar Game Development dan Game Design
  \item Pengenalan Unity Engine 2024+ dan Interface
  \item Dasar Pemrograman C\# untuk Unity (Modern C\# Features)
  \item GameObject, Component, Prefab, dan Scene Management
  \item Unity Input System (New) dan Player Controller
  \item Physics, Collision, dan Rigidbody Simulation
  \item Animation System: Animator Controller dan State Machines
  \item UI/UX Game dengan Unity UI Toolkit
  \item Audio System: Spatial Audio dan Audio Mixer
  \item Level Design dan Game Mechanics Balancing
  \item AI Sederhana: NavMesh dan State Machines
  \item Performance Optimization dan Advanced Debugging
  \item Build Pipeline dan Multi-Platform Deployment
  \item Version Control dengan Git dan Unity Cloud Build
\end{enumerate}

\vspace{0.5cm}

% ============================================================
% 6. Metode Pembelajaran
% ============================================================
\section*{6. Metode Pembelajaran}
\begin{itemize}[leftmargin=*]
  \item \textbf{Flipped Classroom}: Mahasiswa mempelajari materi melalui Unity Learn sebelum sesi kelas, diskusi mendalam di kelas.
  \item \textbf{Project-Based Learning (PBL)}: Pengembangan mini‑game tiap modul dengan milestone yang jelas.
  \item \textbf{Problem-Based Learning}: Penyelesaian studi kasus nyata dari industri game development.
  \item \textbf{Praktikum Terbimbing}: Hands-on coding di Laboratorium Unity dengan bimbingan asisten.
  \item \textbf{Peer Programming dan Code Review}: Kolaborasi tim dan review kode sesama mahasiswa.
  \item \textbf{Game Jam Sessions}: Kompetisi pengembangan game dalam waktu terbatas.
  \item \textbf{Industry Guest Lectures}: Sesi bersama praktisi game development.
  \item \textbf{Portfolio Development}: Pembangunan portfolio game untuk persiapan karir.
\end{itemize}

\vspace{0.5cm}

% ============================================================
% 7. Pengalaman Belajar Mahasiswa
% ============================================================
\section*{7. Pengalaman Belajar Mahasiswa}
\begin{itemize}[leftmargin=*]
  \item Menyusun Game Design Document (GDD) komprehensif untuk tiap proyek.
  \item Mengembangkan prototype game 2D/3D menggunakan Unity 2024+ dengan best practices.
  \item Implementasi kontrol pemain, fisika realistis, animasi fluida, UI responsif, dan audio immersive.
  \item Melakukan iterative development berdasarkan playtesting dan feedback systematic.
  \item Optimasi performa menggunakan Unity Profiler dan advanced optimization techniques.
  \item Menerapkan version control dengan Git dan collaborative development workflows.
  \item Membuat build multi-platform dengan proper settings dan quality assurance.
  \item Menulis technical documentation dan project reports dengan standar industri.
  \item Mempresentasikan game demo dengan professional pitch dan technical explanation.
  \item Membangun portfolio GitHub dengan game projects dan code samples.
\end{itemize}

\vspace{0.5cm}

% ============================================================
% 8. Kriteria, Indikator, dan Bobot Penilaian
% ============================================================
\section*{8. Kriteria, Indikator, dan Bobot Penilaian}
\begin{longtable}{|>{\raggedright\arraybackslash}p{2.5cm}|>{\raggedright\arraybackslash}p{4cm}|>{\raggedright\arraybackslash}p{5.5cm}|c|}
\hline
\textbf{Komponen} & \textbf{Teknik Asesmen} & \textbf{Indikator / CPMK} & \textbf{Bobot (\%)}\\
\hline
\endfirsthead
\hline
\textbf{Komponen} & \textbf{Teknik Asesmen} & \textbf{Indikator / CPMK} & \textbf{Bobot (\%)}\\
\hline
\endhead
\hline
\endfoot

Tugas Individu & Coding Assignment \& GitHub Repository & Sub‑CPMK 1‑4 & 15\\
\hline
Kuis & MCQ \& Essay \& Practical Coding & Sub‑CPMK 5‑7 & 10\\
\hline
Praktikum & Lab Exercise \& Unity Project Demo & Sub‑CPMK 8‑10 & 15\\
\hline
UTS & Written Exam \& Coding Challenge \& Debugging Task & CPMK‑1, CPMK‑2, CPMK‑3 & 20\\
\hline
Proyek Kelompok & Game Development Project \& Portfolio \& Documentation & CPMK‑11, CPMK‑12, CPMK‑13 & 20\\
\hline
UAS & Comprehensive Exam \& Final Game Demo \& Oral Defense & CPMK‑1‑13 & 20\\
\hline
\textbf{Total} &  &  & \textbf{100}\\
\hline
\end{longtable}

\textbf{Kriteria Penilaian Berbasis OBE:}
\begin{itemize}[leftmargin=*]
  \item \textbf{A (85‑100)}: Menguasai semua CPMK dengan sangat baik, mampu menerapkan dalam kasus kompleks, menunjukkan inovasi dan kreativitas.
  \item \textbf{B (70‑84)}: Menguasai sebagian besar CPMK dengan baik, mampu mengaplikasikan konsep secara konsisten.
  \item \textbf{C (60‑69)}: Menguasai CPMK dasar dengan cukup, memerlukan bimbingan minimal untuk kasus kompleks.
  \item \textbf{D (50‑59)}: Menguasai sebagian kecil CPMK, kesulitan dalam aplikasi praktis.
  \item \textbf{E (<50)}: Belum menguasai CPMK yang ditetapkan, perlu remedial intensif.
\end{itemize}

\textbf{Rubrik Penilaian Proyek Game:}
\begin{itemize}[leftmargin=*]
  \item \textbf{Gameplay Mechanics (30\%)}: Kontrol responsif, mekanik jelas, balancing tepat.
  \item \textbf{Technical Implementation (25\%)}: Kode bersih, struktur OOP, performa optimal.
  \item \textbf{Visual \& Audio (20\%)}: Visual menarik, audio sesuai, UI intuitif.
  \item \textbf{Design \& Innovation (15\%)}: Konsep orisinal, level design menarik.
  \item \textbf{Documentation \& Presentation (10\%)}: Dokumentasi lengkap, presentasi jelas.
\end{itemize}

\vspace{0.5cm}

% ============================================================
% 9. Evaluasi dan Refleksi Pembelajaran (opsional)
% ============================================================
\section*{9. Evaluasi dan Refleksi Pembelajaran}
\begin{itemize}[leftmargin=*]
  \item \textbf{Evaluasi Formatif}: Kuis mingguan, tugas coding, peer review, serta feedback berkelanjutan.
  \item \textbf{Evaluasi Sumatif}: UTS, UAS, dan penilaian proyek akhir.
  \item \textbf{Refleksi Mahasiswa}: Jurnal belajar mingguan dan self‑assessment setelah tiap milestone.
  \item \textbf{Evaluasi Dosen}: Survey kepuasan mahasiswa pada pertengahan dan akhir semester.
  \item \textbf{Continuous Improvement}: Analisis hasil penilaian untuk perbaikan RPS pada semester berikutnya.
\end{itemize}

\vspace{0.5cm}

% ============================================================
% 10. Daftar Referensi
% ============================================================
\section*{10. Daftar Referensi}
\begin{enumerate}[leftmargin=*]
  \item Unity Technologies. \textit{Unity User Manual 2024}. \url{https://docs.unity3d.com/Manual/}
  \item Unity Learn. \textit{Junior Programmer Pathway}. 2024. \url{https://learn.unity.com/pathway/junior-programmer}
  \item Gregory, J. (2018). \textit{Game Engine Architecture} (3rd ed.). A K Peters/CRC Press.
  \item Gregory, J. (2019). \textit{Game Coding Complete} (4th ed.). Mercury Learning and Information.
  \item Nystrom, R. (2014). \textit{Game Programming Patterns}. Genever Benning.
  \item Bishop, L. (2023). \textit{Unity in Action: Multiplatform Game Development in C\#} (3rd ed.). Manning Publications.
  \item Thorn, A. (2022). \textit{Mastering Unity 2D Game Development}. Packt Publishing.
  \item Sharp, J. (2023). \textit{C\# in Depth} (4th ed.). Manning Publications.
  \item Microsoft. \textit{C\# Programming Guide}. 2024. \url{https://learn.microsoft.com/en-us/dotnet/csharp/}
  \item Microsoft. \textit{Unity Scripting API Reference}. 2024. \url{https://docs.unity3d.com/ScriptReference/}
  \item Gregory, J. (2020). \textit{3D Math Primer for Graphics and Game Development} (2nd ed.). A K Peters/CRC Press.
  \item Buckland, M. (2022). \textit{Programming Game AI by Example} (2nd ed.). CRC Press.
  \item Lengyel, E. (2022). \textit{Mathematics for 3D Game Programming and Computer Graphics} (4th ed.). A K Peters/CRC Press.
  \item Gamma, E., Helm, R., Johnson, R., \& Vlissides, J. (1994). \textit{Design Patterns: Elements of Reusable Object‑Oriented Software}. Addison‑Wesley.
  \item Martin, R. C. (2017). \textit{Clean Architecture: A Craftsman's Guide to Software Structure and Design}. Prentice Hall.
  \item Freeman, E., \& Robson, E. (2020). \textit{Head First Design Patterns} (2nd ed.). O'Reilly Media.
  \item Fowler, M. (2018). \textit{Refactoring: Improving the Design of Existing Code} (2nd ed.). Addison‑Wesley Professional.
  \item GitHub. \textit{Git Documentation}. 2024. \url{https://git-scm.com/doc}
  \item Unity Technologies. \textit{Unity Profiler Documentation}. 2024. \url{https://docs.unity3d.com/Manual/Profiler.html}
\end{enumerate}

\vspace{0.5cm}

% ============================================================
% 11. Sumber Belajar dan Platform Digital
% ============================================================
\section*{11. Sumber Belajar dan Platform Digital}

\textbf{Platform Pembelajaran Utama:}
\begin{itemize}[leftmargin=*]
  \item \textbf{Unity Learn}: \url{https://learn.unity.com/} - Platform resmi Unity dengan tutorial dan pathway
  \item \textbf{Unity Documentation}: \url{https://docs.unity3d.com/} - Dokumentasi teknis lengkap
  \item \textbf{GitHub Classroom}: Untuk submission tugas dan version control
  \item \textbf{Discord Server}: Komunikasi, Q\&A, dan kolaborasi
\end{itemize}

\textbf{Software dan Tools yang Dibutuhkan:}
\begin{itemize}[leftmargin=*]
  \item Unity Editor 2024.2+ (Personal Edition - Gratis)
  \item Visual Studio 2022 Community atau Visual Studio Code
  \item Git for Windows dan GitHub Desktop
  \item Unity Hub untuk mengelola multiple Unity versions
\end{itemize}

\textbf{Hardware Minimum Requirements:}
\begin{itemize}[leftmargin=*]
  \item OS: Windows 10/11 64-bit, macOS 10.15+, atau Linux modern
  \item CPU: Quad-core 2.5 GHz atau lebih tinggi
  \item RAM: 8 GB minimum (16 GB direkomendasikan)
  \item GPU: DirectX 11 compatible dengan 2 GB VRAM
  \item Storage: 10 GB free space untuk Unity dan project files
\end{itemize}

\textbf{Online Learning Resources:}
\begin{itemize}[leftmargin=*]
  \item Unity Learn Junior Programmer Pathway
  \item Unity Creator Portal (tutorials, case studies)
  \item YouTube Unity Official Channel
  \item Game Development Stack Exchange
  \item Unity Asset Store (free assets untuk prototyping)
\end{itemize}

\vspace{0.5cm}

% ============================================================
% 12. Kebijakan Akademik dan Etika
% ============================================================
\section*{12. Kebijakan Akademik dan Etika}

\textbf{Kebijakan Kehadiran:}
\begin{itemize}[leftmargin=*]
  \item Minimal kehadiran 75\% untuk dapat mengikuti UAS
  \item Kehadiran praktikum wajib 100\% kecuali dengan surat dokter
  \item Toleransi keterlambatan 15 menit
\end{itemize}

\textbf{Kebijakan Tugas:}
\begin{itemize}[leftmargin=*]
  \item Tugas dikumpulkan melalui GitHub Classroom sesuai deadline
  \item Late submission: pengurangan 10\% per hari (maksimal 3 hari)
  \item Plagiarisme akan mendapat nilai 0 dan konsekuensi akademik
\end{itemize}

\textbf{Code of Conduct:}
\begin{itemize}[leftmargin=*]
  \item Menghargai kolaborasi dan peer review
  \item Menggunakan version control dengan proper commit messages
  \item Mengikuti best practices dalam code documentation
  \item Menjaga profesionalisme dalam semua komunikasi
\end{itemize}

\vspace{1cm}
\begin{flushright}
\begin{tabular}{c}
Disusun oleh,\\[2cm]
\textbf{[Nama Dosen], M.Kom.}\\
NIP. [Nomor Induk Pegawai]\\[0.5cm]
\textbf{Menyetujui,}\\[2cm]
\textbf{Ketua Program Studi}\\
\textbf{Teknik Informatika}\\
\textbf{[Nama Ketua Prodi], M.Kom., Ph.D.}\\
NIP. [Nomor Induk Pegawai]
\end{tabular}
\end{flushright}

\vspace{0.5cm}
\begin{center}
\small\textit{Silabus ini disusun berdasarkan Kurikulum Berbasis OBE dan mengacu pada standar industri game development terkini (2024)}
\end{center}

\end{document}
